%============================================================
% CH1 — SOLUTION SKETCHES FOR THE PRACTICE SET   (tag g1s)
% Graph model, algorithm and complexity for each problem.
% Sketches, not full code: the point is to check your MODEL.
% Every sketch was implemented and reproduces the official samples.
% Order and groups follow problem_set.tex (Tiers 1-3, then Stretch).
%============================================================
\subsection{Solution Sketches}
\label{ssec:ch1-sketches}

\begin{warningbox}[Read a sketch only after a real attempt]
Each sketch gives the \emph{design} --- the graph model, the algorithm, the complexity and the one trap of the problem --- in a few lines, so you can compare it with your own Graph Model. Writing and debugging the code is left to you. If your model differs but you can argue why it answers the question, it is probably also right.
\end{warningbox}

\begingroup
\providecommand{\sk}{}
\renewcommand{\sk}[2]{\par\smallskip\noindent\textbf{\textcolor{deepblue}{#1}}\quad #2\par}

\begin{solutionbox}[Tier 1 --- Warm-up]
\sk{P1.01 Counting Roads (ABC061 B, difficulty 108).}{For each road $(a,b)$: \texttt{deg[a]++}, \texttt{deg[b]++}; print \texttt{deg[1..N]}. Parallel roads are separate edges (sample~2 prints $5$ and $5$). \bigO{N+M}.}
\sk{P1.02 Adjacency Matrix (ABC343 B, difficulty $<0$).}{Row $i$ of the matrix is vertex $i$'s neighbour set: print every $j$ with $A_{ij} = 1$, in increasing order (an empty line for an isolated vertex, sample~2). This is exactly the conversion from a matrix to adjacency lists, \bigO{N^2} --- the input size.}
\sk{P1.03 Path Graph? (ABC287 C, difficulty 466).}{Yes iff $M = N-1$, every degree is $\le 2$, and one DFS from vertex $1$ reaches all $N$ vertices. \emph{Trap:} all three tests are needed --- a triangle plus an isolated vertex ($N = 4$, $M = 3$) passes the first two and fails the third. \bigO{N+M}.}
\sk{P1.04 Love Triangle (CF 939A).}{Functional graph $i \to f_i$ with $f_i \ne i$. A love triangle is a cycle of length exactly $3$; since $f_i \ne i$, the test $f(f(f(i))) = i$ for some $i$ cannot be fooled by a cycle of length $1$. \bigO{n}.}
\sk{P1.05 Computer Game (CF 1598A).}{Vertices: the safe cells of the $2 \times n$ grid; edges: the 8 neighbour moves. With diagonal moves, two consecutive columns that each have a safe cell are always connected; the path is cut only by a column whose two cells are both traps. Answer \texttt{NO} iff such a column exists (a BFS on the $2n$ cells gives the same answer, checked on random grids). \bigO{n}.}
\sk{P1.06 Li Hua and Maze (CF 1797A).}{Grid graph; $\deg(\text{cell}) = $ number of in-grid side neighbours ($2$, $3$ or $4$). Blocking all neighbours of the start (or of the end) cell disconnects them, so the answer is $\le \min(\deg s, \deg t)$; because $n, m \ge 4$ and the cells are not adjacent, there are that many vertex-disjoint routes, so fewer blocks never suffice. \emph{Trap:} the tempting BFS is impossible ($n, m \le 10^9$) and unnecessary. Answer $\min(\deg s, \deg t)$ in \bigO{1} --- checked against an exhaustive search on all $4 \times 4$ to $5 \times 5$ grids.}
\sk{P1.07 Party (CF 115A).}{Arrow from each employee to his manager ($-1$ for a root): a forest of parent pointers. Employees at the same depth are never superior to one another, and a chain of depth $d$ needs $d$ groups, so the answer is the maximum depth (counting the root as $1$). Walk up from every vertex (\bigO{n^2} with $n \le 2000$) or memoise depths for \bigO{n}. Sample: $3$.}
\sk{P1.08 Begginer's Zelda (CF 1905B).}{One operation lowers the number of leaves by at most two (only the ends of the path can stop being leaves), and contracting the path between two leaves lowers it by exactly two until at most two leaves are left. So with $L$ leaves the answer is $\lceil L/2 \rceil$. Count degree-$1$ vertices: \bigO{n}. Samples: $1, 3, 2, 2$; checked against an exhaustive search over all contraction sequences for $n \le 7$.}
\end{solutionbox}

\begin{solutionbox}[Tier 2 --- Core]
\sk{P1.09 Validate Binary Tree Nodes.}{Arrows from each node to its children. A valid binary tree needs: every in-degree $\le 1$; exactly one vertex with in-degree $0$ (the root); and the root reaches all $n$ vertices. \emph{Trap:} the degree tests alone accept a tree plus a separate cycle such as $2 \to 3 \to 2$, whose vertices all have in-degree $1$. One counting pass and one DFS. \bigO{n}. Samples: \texttt{true}, \texttt{false}, \texttt{false}.}
\sk{P1.10 Coloring A Border.}{Flood fill the component of \texttt{(row, col)} (same colour, 4 neighbours). A component cell is a \emph{border} cell if it lies on the grid edge or has a side neighbour of a different colour. \emph{Trap:} collect the border cells during the search and recolour them afterwards --- recolouring during the search changes the colours that later cells compare against. \bigO{mn}. All three samples reproduced.}
\sk{P1.11 Remove Methods From Project.}{Arrow $a \to b$: ``$a$ invokes $b$''. The suspicious set $S$ is everything reachable from $k$ (one DFS). Then scan all invocations once: if some $a \notin S$ invokes $b \in S$, nothing may be removed and the answer is all methods; otherwise the answer is the methods outside $S$. \bigO{n+m}. Samples: $[0,1,2,3]$, $[3,4]$, $[\,]$.}
\sk{P1.12 Martial artist (ABC226 C, difficulty 539).}{Arrow $i \to A_{i,j}$ (``move $i$ needs move $A_{i,j}$''). The moves that must be learned are exactly those reachable from $N$; learning them in increasing index order respects every prerequisite ($A_{i,j} < i$). Answer: $\sum T_i$ over the reachable set, in \texttt{long long} (up to $2 \cdot 10^{14}$). \bigO{N + \sum K_i}. Samples: $10$, $5000000000$.}
\sk{P1.13 Ladder Takahashi (ABC277 C, difficulty 540).}{Vertices: floor $1$ and every floor touched by a ladder (at most $2N+1$ of them); undirected edge per ladder. Labels go up to $10^9$, so use \texttt{map<int, vector<int>>} (or compress coordinates). DFS from floor $1$; answer the largest visited floor ($1$ if no ladder touches floor $1$). \bigO{N \log N}. Samples: $10$, $12$, $1$.}
\sk{P1.14 Little Pigs and Wolves (CF 116B).}{Bipartite graph wolves--pigs on the grid. Each pig has at most one wolf neighbour, so different wolves never compete for a pig: every wolf with at least one pig neighbour eats one. Answer: the number of wolves with a pig among their 4 neighbours. \bigO{nm}. Samples: $2$, $0$.}
\sk{P1.15 Students and Shoelaces (CF 129B).}{Repeat: collect \emph{all} vertices of degree $1$ first, then delete them and their edges; count the rounds until no vertex has degree $1$. Deleting while scanning would let a vertex that only became a leaf in this round leave too early. At most $n$ rounds of \bigO{n+m} each, fine for $n \le 100$. Samples: $0$, $2$, $1$.}
\sk{P1.16 Mocha and Hiking (CF 1559C).}{Graph: spine $1 \to 2 \to \dots \to n$ and village $n{+}1$ joined to each $i$ by $i \to n{+}1$ ($a_i = 0$) or $n{+}1 \to i$ ($a_i = 1$). If $a_1 = 1$: $n{+}1, 1, \dots, n$. Else if $a_n = 0$: $1, \dots, n, n{+}1$. Else $a$ starts with $0$ and ends with $1$, so some $i$ has $a_i = 0$, $a_{i+1} = 1$: $1, \dots, i, n{+}1, i{+}1, \dots, n$. \emph{Trap:} the statement allows $-1$, but a route always exists --- searching for a Hamiltonian path is never needed. \bigO{n}. Any valid route is accepted; ours was checked by a route validator on the samples and on random inputs.}
\sk{P1.17 Arrow Path (CF 1948C).}{State: the cell where the robot stands after a full second (move, then forced arrow step); start $(1,1)$. From a state, try the 4 neighbours inside the grid and apply the arrow of that neighbour (if it stays inside). BFS over the $\le 2n$ states; answer \texttt{YES} iff $(2,n)$ is reached. \bigO{n}. Samples: \texttt{YES YES NO YES}; counting $(2,n)$ also when it is only touched before the arrow step gives the same answers on random grids.}
\end{solutionbox}

\begin{solutionbox}[Tier 3 --- Challenge]
\sk{P1.18 Path Existence Queries in a Graph I.}{\texttt{nums} is sorted, and an edge joins $i$, $j$ iff $|nums_i - nums_j| \le \texttt{maxDiff}$. If $i < j$ are connected by any path, every gap between consecutive values in $[i, j]$ is $\le \texttt{maxDiff}$, and then consecutive indices are adjacent. So label components left to right: $c_0 = 0$, $c_i = c_{i-1} + [nums_i - nums_{i-1} > \texttt{maxDiff}]$, and answer each query by $c_u = c_v$. The graph may have $\Theta(n^2)$ edges --- never build it. \bigO{n + q}. Samples reproduced; checked against BFS on the explicit graph.}
\sk{P1.19 Grid and Magnet (ABC351 D, difficulty 974).}{Free cells are vertices. A free cell next to a magnet is a \emph{dead end}: it can be entered but not left. Flood fill from every free non-dead-end cell not yet in a region, expanding only non-dead-end cells, and count the cells reached. \emph{Trap:} one dead end can touch several regions, so mark visits with a region stamp, not with one global \texttt{visited} flag; and a lonely dead-end cell gives freedom $1$. \bigO{HW}. Samples: $9$, $1$; checked against a BFS from every cell.}
\sk{P1.20 Road Construction (CF 330B).}{Any valid network is connected, so it needs $\ge n-1$ roads. The $m < n/2$ forbidden pairs touch at most $2m < n$ cities, so some city $c$ appears in none of them. The star with centre $c$ uses $n-1$ allowed roads and any two cities are at distance $\le 2$. \bigO{n+m}. Sample: $3$ roads (a different star is also accepted).}
\sk{P1.21 Permutation Game (CF 1968D).}{A player at $P$ visits $v_0 = P, v_1 = p_{v_0}, \dots$ (a cycle of the permutation). Optimal play: move for $i$ turns, then stay on $v_i$ for the remaining $k - i$ turns: $\text{score}_i = \sum_{j<i} a_{v_j} + (k-i)\, a_{v_i}$. Positions repeat after at most $n$ moves, so take $\max$ over $i < \min(n, k)$, in \texttt{long long} (up to $10^{18}$). Compare the two players. \bigO{n} per test. All $10$ sample verdicts reproduced; also checked against trying all $2^k$ stay/move choices on small inputs.}
\sk{P1.22 Find the Permutation (CF 2056B).}{For values $a < b$, the edge $a$--$b$ exists iff $a$ stands before $b$ in $p$; without the edge, $b$ stands before $a$. So the matrix decides the order of every pair: sort the values $1..n$ with this comparator (it is consistent because it comes from a real permutation). \bigO{n^2}, the input size. Samples: $1$; $4\,2\,1\,3\,5$; $6\,5\,4\,3\,2\,1$; checked by building the graph of random permutations and recovering them.}
\sk{P1.23 Divide by three, multiply by two (CF 977D).}{Vertices: the numbers; arrows $x \to x/3$ and $x \to 2x$ when present. No number can have both successors in the list (reaching the other one later would need $2^a = 3^b$ with $a + b > 0$), so the graph is a single path: start at the number with in-degree $0$ and follow the arrows. \bigO{n^2} with $n \le 100$, or a \texttt{set} for \bigO{n \log n}; values reach $3 \cdot 10^{18}$, so \texttt{long long}. Equivalent shortcut: sort by (more factors $3$ first, then smaller value). Samples reproduced.}
\sk{P1.24 Ball in Berland (CF 1475C).}{Bipartite graph boys--girls with the $k$ pairs as edges. The edges that share an endpoint with edge $(a, b)$, itself included, number $\deg(a) + \deg(b) - 1$; the remaining $k - \deg(a) - \deg(b) + 1$ are compatible. Summing over all edges counts every compatible pair twice: answer $= \frac12 \sum_e (k - \deg a - \deg b + 1)$, in \texttt{long long}. \bigO{a + b + k}. Samples: $4$, $0$, $2$; checked against all pairs on random inputs.}
\sk{P1.25 Web of Lies (CF 1548A).}{Claim: the survivors are exactly the nobles with no friend of higher power. Such a noble is never vulnerable (his friends only disappear). A noble $v$ with a stronger friend $w$ dies: $w$ cannot die while $v$ is alive (then $w$ has a weaker living friend), and as long as some living noble has a stronger living friend, the weakest such noble has only stronger living friends, so the process does not stop. \emph{Trap:} simulating the process per query is far too slow. Maintain $\text{cnt}[v]$ = number of stronger friends: on adding $(u, v)$ increment $\text{cnt}[\min(u,v)]$, on removing decrement it, and keep the number of $v$ with $\text{cnt}[v] = 0$ up to date. \bigO{1} per query. Samples: $2, 1$ and $1$; checked against simulating the process on random small inputs.}
\sk{P1.26 Trapped in the Witch's Labyrinth (CF 2034C).}{Compute the set $E$ of cells that \emph{must} escape, by a reverse search from the border: a fixed arrow is in $E$ if it points outside or into $E$; a \texttt{?} is in $E$ only if all four neighbours are outside or in $E$. Every other cell can be trapped: point a \texttt{?} at a neighbour not in $E$ (two neighbouring \texttt{?} cells point at each other). Answer $nm - |E|$. \emph{Trap:} a lonely \texttt{?} surrounded by escaping cells cannot be saved. \bigO{nm}. Samples: $0$, $6$, $5$; checked against trying every assignment of the \texttt{?} cells on small grids.}
\end{solutionbox}

\begin{solutionbox}[Stretch --- above this chapter]
\sk{P1.27 Reachability in Functional Graph (ABC357 E, difficulty 1295).}{Each component is a cycle with trees hanging into it. A cycle vertex reaches exactly the $L$ vertices of its cycle; a tail vertex $u$ reaches itself plus everything $a_u$ reaches: $r(u) = 1 + r(a_u)$. Find the cycles with colours $0/1/2$ (a walk that meets a vertex of colour $1$ has closed a new cycle), then fill the tails by walking until a known $r$ and assigning backwards. Answer $\sum_u r(u)$, up to $4 \cdot 10^{10}$: \texttt{long long}. \bigO{N}. Samples: $8$, $14$, $41$; checked against brute force.}
\sk{P1.28 Lunar New Year and a Wander (CF 1106D).}{Bob may walk back through visited vertices, so at any moment he can record \emph{any} discovered but unrecorded vertex. Greedy: keep the frontier in a min-heap, always record the smallest, then push its undiscovered neighbours (marked on push, as in BFS). Exchange argument: recording anything else first makes the sequence lexicographically larger. \bigO{m \log n}. Samples reproduced.}
\end{solutionbox}

%------------------------------------------------------------
\subsection{Answers to the Self-Check Quiz}
\label{ssec:ch1-quiz-answers}

\begin{solutionbox}[Quiz answers]
\begin{enumerate}[leftmargin=*]
    \item No. The degrees sum to $11$, but every graph has $\sum_v \deg(v) = 2|E|$, an even number (equivalently: it would have three odd-degree vertices, an odd number).
    \item $2m$ entries: every undirected edge appears in the lists of both endpoints.
    \item The matrix has $n^2 = 4 \cdot 10^{10}$ cells --- tens of gigabytes even at one byte per cell, and \bigO{n^2} time just to clear it. Use adjacency lists: \bigO{n+m} memory.
    \item Thanks to \texttt{visited}, every vertex is pushed and processed at most once, and processing $v$ scans its list once; the scans add up to $\sum_v \deg(v) = 2m$ (or $m$ for a directed graph). Total \bigO{n+m}.
    \item Marking on pop lets the same vertex be pushed once per neighbour that sees it before it is popped, so the queue can hold \bigO{m} copies and work is repeated. Marking on push guarantees each vertex enters the queue exactly once (and, for distances, fixes its distance at the first, smallest value).
    \item Every recursive call takes a stack frame; $10^6$ nested calls exceed the default stack (a few megabytes), so the program dies with a stack overflow. Use an iterative DFS with an explicit \texttt{stack} (or BFS), or make sure the depth is small.
    \item With $n-1$ edges: a tree. With exactly $n$ edges: a tree plus one edge, i.e.\ a connected graph with exactly one cycle.
    \item The judge is defined by degrees only: trusted by all $n-1$ others (in-degree $n-1$) and trusting nobody (out-degree $0$). At most one vertex can satisfy both (two judges would have to trust each other), so one counting pass and one check per vertex suffice.
    \item $RC$ vertices; $R(C-1)$ horizontal plus $C(R-1)$ vertical edges, i.e.\ $2RC - R - C$ edges --- fewer than $2RC$.
    \item $10^4$ states (\texttt{0000} to \texttt{9999}); each has $8$ neighbours (4 wheels, each turned up or down).
    \item From a value $y \ge m$ at least $y - m$ presses are needed (a press lowers the value by at most $1$), and $y - m$ presses of the ``$-1$'' button achieve it, so an optimal sequence never doubles a value $\ge m$. Doubling therefore only happens to values $< m$, giving values $< 2m$, and otherwise values only go down from $n$. Hence every value on an optimal path is $\le \max(n, 2m)$.
    \item One directed cycle with trees hanging into it, all arrows of the trees pointing towards the cycle (the ``$\rho$'' shape). Every vertex has an outgoing arrow, so a walk never stops; with only $n$ vertices, by the pigeonhole principle some vertex repeats within $n+1$ steps, and from then on the walk goes round the cycle forever.
\end{enumerate}
\end{solutionbox}
\endgroup
